\documentclass[10pt,a4paper]{article}
\usepackage{amsmath,amssymb,amsthm,graphicx}
\usepackage{amsmath,amssymb,amsthm,hyperref,xfrac}
\providecommand{\ie}{i.e.\ }
\providecommand{\skpt}{\smallskip}
\usepackage[scr=rsfs,cal=euler]{mathalfa}
\multlinegap0pt

\hyphenation{begin decom-po-si-tion sequence}
\newenvironment{enumeratea}
{\bgroup\def\theenumi{\alph{enumi}}\begin{enumerate}}
{\end{enumerate}\egroup}

\newcommand{\ba}{{\boldsymbol a}}
\newcommand{\bb}{{\boldsymbol b}}
\newcommand{\bd}{{\boldsymbol d}}
\newcommand{\bu}{{\boldsymbol u}}
\newcommand{\bx}{{\boldsymbol x}}
\newcommand{\bmy}{{\boldsymbol y}}

\newcommand{\bD}{{\boldsymbol D}}
\newcommand{\bL}{{\boldsymbol L}}
\newcommand{\bP}{{\boldsymbol P}}
\newcommand{\bV}{{\boldsymbol V}}
\newcommand{\bX}{{\boldsymbol X}}
\newcommand{\bZ}{{\boldsymbol Z}}

\newcommand{\Psfrac}[2]{(\sfrac{#1}{#2})}
\newcommand{\psfrac}[2]{\sfrac{(#1)}{#2}}
\newcommand{\bigO}[1]{O(#1)}

\usepackage{mathtools}
\usepackage{bbm}

\newcommand{\NN}{\mathbb{N}}
\newcommand{\ZZ}{\mathbb{Z}}
\newcommand{\RR}{\mathbb{R}}
\newcommand{\CC}{\mathbb{C}}
\newcommand{\PP}{\mathbb{P}}
\newcommand{\EE}{\mathbb{E}}

\newcommand{\boldalpha}{\boldsymbol{\alpha}}
\newcommand{\bolddelta}{\boldsymbol{\delta}}
\newcommand{\boldrho}{\boldsymbol{\rho}}
\newcommand{\dee}{\,\mathrm{d}}
\newcommand{\one}{\mathbbm{1}}
\newcommand{\BC}{\mathcal{B}}
\newcommand{\EC}{\mathcal{E}}
\newcommand{\RS}{\mathscr{R}}
\newcommand{\prob}[1]{\PP[ #1 ]}
\newcommand{\bigoo}[1]{O\!\left(#1 \right)}
\newcommand{\abs}[1]{\left\lvert #1 \right\rvert}
\newcommand{\floor}[1]{\left\lfloor #1 \right\rfloor}
\newcommand{\riemint}[4]{\int_{#1}^{#2}\! #3 \, \mathrm{d}#4}
\newcommand{\maxnorm}[1]{\left\| #1 \right\|_{\infty}}

\renewcommand{\Re}{\textup{Re}}
\renewcommand{\Im}{\textup{Im}}

\newcommand{\dtv}{\mathrm{d}_{\mathrm{TV}}}
\newcommand{\dw}{\mathrm{d}_{\mathrm W}}

\let\epsilon\varepsilon

\newcommand{\pextra}{P_{\mathtt{extra}}}
\newcommand{\parratia}{P_{\mathtt{Arratia}}}

\newcommand{\simplex}{\Delta^{k-1}}

\newcommand{\bg}{\big}
\newcommand{\bgg}{\Big}
\newcommand{\bggg}{\bigg}

\newcommand{\Rnt}{R_{\mathrm{NT}}}
\newcommand{\Rprob}{R_{\mathrm{PD}}}
\newcommand{\ECnt}{\EC_{\mathrm{NT}}}
\newcommand{\ECprob}{\EC_{\mathrm{PD}}}

\newcommand{\Bnt}{\mathcal{B}_{\mathrm{NT}}}
\newcommand{\Bprob}{\mathcal{B}_{\mathrm{PD}}}

\DeclareMathOperator{\vard}{Var}
\DeclareMathOperator{\Dir}{Dir}

\theoremstyle{plain}
\newtheorem{cor}{Corollary}[section]
\newtheorem{lem}[cor]{Lemma}
\newtheorem{prop}[cor]{Proposition}
\newtheorem{conj}[cor]{Conjecture}
\newtheorem{theo}{Theorem}

\theoremstyle{definition}
\newtheorem{example}{Example}
\newtheorem{deff}[cor]{Definition}
\newtheorem{rem}{Remark}[section]
\newtheorem*{rems}{Remark}
\newtheorem*{remss}{Remarks}

\makeatletter
\ams@newcommand{\multiint}[1]{\DOTSI\protect\MultiIntegral{#1}}
\renewcommand{\MultiIntegral}[1]{%
\edef\ints@c{\noexpand\intop
\ifnum#1=\z@\noexpand\intdots@\else\noexpand\intkern@\fi
\replicate{#1-2}{\noexpand\intop\noexpand\intkern@}%
\noexpand\intop
\noexpand\ilimits@
}%
\futurelet\@let@token\ints@a
}
\makeatother
\ExplSyntaxOn
\cs_new:Npn \replicate #1 #2 { \prg_replicate:nn { #1 } { #2 } }
\ExplSyntaxOff
\title{The law of Arratia's random integer J\_x\textasciicircum{}*: P[J\_x\textasciicircum{}* = j] = (1/(j log x))(1 + O(1/log x))}
\author{Tony Haddad, Dimitris Koukoulopoulos}
\date{doi:10.5802/jep.317; CC BY 4.0}
\begin{document}
\maketitle
\noindent\textit{Editorial note.} The claim is Arratia's Lemma 2 as restated by the authors, who supply their own proof at 773-801; the attribution stays in the lemma heading.\par\medskip
\noindent\textit{Editorial note.} The Mecke equation is quoted from Last-Penrose at 1845-1852 and is applied to a nonnegative summand, exactly the range of the quoted statement.\par\medskip
\noindent\textit{Editorial note.} The following section 803-1113 on the event \{J\_x != J\_x\textasciicircum{}*\} and the total variation distance is not included and is not used by this proof.\par\medskip
\noindent\textit{Editorial note.} All displayed estimates are taken from the fixed publisher source unchanged.\par\medskip
\noindent\textit{Allowed standard prerequisites (editorial).}\begin{itemize}
\item Poisson process calculus on a product space. R has intensity exp(-wy) dw dy; with S\_i = sum\_\{l>=i\} Y\_l the process \{S\_i\} is Poisson of intensity ds/s; the Mecke equation evaluates sums over distinct Poisson points.. Supplies R\textasciicircum{}*, its labelling, and the Mecke identity that converts the indicator of \{J\_x\textasciicircum{}* = j\} into an integral. The spacing lemma is quoted, not reproved.
\item Chebyshev function and Mertens estimate. theta(x) = sum\_\{p<=x\} log p, Lambda is von Mangoldt, and sum\_\{q<=u\} Lambda(q)/q = log u + O(1).. Turns the prime-power tail in the Mecke integral into an explicit asymptotic in t.
\item Dirichlet series of a multiplicative Poisson product. For N\_q(t) Poisson of parameter Lambda(q)/(q\textasciicircum{}\{1+t\} log q) and I\_t\textasciicircum{}* = prod\_q q\textasciicircum{}\{N\_q(t)\}, one has E[(I\_t\textasciicircum{}*)\textasciicircum{}\{-s\}] = zeta(1+t+s)/zeta(1+t).. Identifies P[I\_t\textasciicircum{}* = i] = i\textasciicircum{}\{-1-t\}/zeta(1+t).
\end{itemize}\medskip
There are many ways to define the Poisson--Dirichlet distribution. One of the most intuitive ones involves a ``stick-breaking'' process that we will use throughout the paper. We~start by sampling a sequence of i.i.d.~random variables $(U_i)_{i \ge 1}$ that are all uniformly distributed in $[0, 1]$. We~then define the sequence $(L_i)_{i \ge 1}$ in the following way:
\[
L_1 \coloneqq U_1 \quad \text{and}\quad L_j \coloneqq U_j \prod_{i=1}^{j-1} (1-U_i) \quad \text{for $j \ge 2$.}
\]
The distribution of the process $\bL = (L_1, L_2, \ldots)$ is called the \textit{GEM distribution} (of~parameter 1), named after Griffiths, Engen and McCloskey (see \cite[Ch.\,41]{JKB} for a discussion on the history of this distribution). Lastly, we~sort the components of $\bL$ in non-increasing order to create $\bV = (V_1, V_2, \ldots)$. The distribution of this process is the \textit{Poisson--Dirichlet distribution} (of parameter 1).\footnote{The GEM and Poisson--Dirichlet distributions have more general definitions involving typically a parameter $\theta$. In~the rest of the paper, we~will not be mentioning the parameter since we will always work with $\theta = 1$.} We note that both $\sum_{i \ge 1} L_i$ and $\sum_{i \ge 1} V_i$ are equal to $1$ almost surely.

\begin{deff}[The Poisson process $\RS$]\label{dfn:Poisson Process R}
\begin{enumeratea}
\item We denote by $\RS$ the Poisson process on $\RR_{>0}^2$ that has intensity measure $e^{-wy}\, \mathrm dw\, \mathrm dy$. We~may assume that the $w$-coordinates of the points of $\RS$ are all distinct since this happens almost surely.
\item We index the points of $\RS = \bg\{(W_i, Y_i) : i \in\ZZ\bg\}$ according to the following rules:
\begin{itemize}
\item $W_i < W_{i+1}$ for all $i\in\ZZ$;
\item if we let $S_i\coloneqq \sum_{\ell \ge i} Y_\ell$ for all $i\in\ZZ$, then we have $S_1 \le \log x < S_0$.
\end{itemize}
\end{enumeratea}
\end{deff}

\begin{rems}
By the mapping theorem (Proposition \ref{prel:map}), projecting $\RS$ on the $w$-axis yields a Poisson Process with intensity $\sfrac{\mathrm dw}{w}$ and, similarly, projecting $\RS$ on its $y$\nobreakdash-axis yields a Poisson Process with intensity $\sfrac{\mathrm dy}{y}$. Therefore, the $w$-coordinates of the points in $\RS$ have almost surely exactly one limit point at $0$ and they are almost surely unbounded. Hence, the indexing $(W_i,Y_i)$ in part (b) of the above definition is well-defined.
\end{rems}

The following lemma describes the distribution of the point process $S_i$.

\begin{lem}[Scale-invariant spacing lemma]
\label{lem:scale-invariant-spacing-lemma}
The point process $\{S_i : i \in \ZZ\}$ is a Poisson process on the positive real line with intensity measure $\sfrac{\mathrm ds}{s}$.
\end{lem}

\section{An integer-friendly version of \texorpdfstring{$J_x$}{Jx}}
\label{sec:integer friendly J}

Let $\Lambda$ be the von Mangoldt function, that is to say
\[
\Lambda(n) = \begin{cases}
\log p &\text{if $n = p^k$ for some prime power $p^k$,} \\
0 &\text{otherwise.}
\end{cases}
\]
Recall the Poisson Process $\RS$ given in Definition \ref{dfn:Poisson Process R}. We~then define
\[
\RS^* \coloneqq \bg\{ \bg( WY/ h(Y), e^{h(Y)} \bg): (W,Y)\in \RS,\ Y > e^{-\gamma}\bg\}.
\]
Without loss of generality, we~may assume that the points $WY/h(Y)$ with $(W,Y)\in\RS$ and $Y>e^{-\gamma}$ are all distinct. By the mapping theorem (Proposition \ref{prel:map}) applied to the map $(w,y)\mapsto (wy/h(y),e^{h(y)})$ on the Poisson process $\RS$ restricted to $\RR_{>0} \times \RR_{>e^{-\gamma}}$, the random set $\RS^*$ is a Poisson process on the space $\RR_{>0}\times \{\text{prime powers}\}$ with mean measure $\mu^*$ satisfying
\[
\mu^*(B \times \{q\}) = \int_{B} \frac{\Lambda(q)}{q^{1+t}} \, \mathrm dt
\]
for any $B\subseteq\RR_{>0}$ and any $q\in\NN$. Note that the $w$-coordinates of the points in~$\RS$ restricted to $\RR_{>0} \times \RR_{>e^{-\gamma}}$ are now bounded with probability one. For each such realization of the points of $\RS^*$, there is a unique labeling of them as \hbox{$\{(T^*_i, Q^*_i):i \in \ZZ_{\le K}\}$} in such a way that the following properties hold:
\begin{itemize}
\item $T^*_{i-1} <T^*_{i}$ for all $i\in\ZZ_{\le K}$;
\item $\prod_{i=1}^K Q_i^*\le x< \prod_{i=0}^K Q_i^*$.
\end{itemize}
(Note that $K$ is a random variable and is not fixed.) We then define the random integer
\[
J_x^*\coloneqq \prod_{i=1}^K Q^*_i.
\]

\subsection{Poisson Processes}

The following definition and propositions are borrowed from Kingman's book on Poisson processes \cite{Kingman93}. Let $(S, \mathcal S)$ be a measurable space with~$S$ being a subset of $\RR^d$ for some $d \ge 1$. A \textit{Poisson process} on a state space $S$ with mean measure $\mu$ is a random countable subset $\Pi \subseteq S$ such that:
\begin{itemize}
\item for any disjoint measurable subsets $A_1, \ldots, A_n$ of $S$, the random variables \hbox{$\#(\Pi \cap A_1)$, \dots, $\#(\Pi \cap A_n)$} are independent;

\item the random variable $\#(\Pi \cap A)$ is a Poisson random variable of parameter $\mu(A)$ for any $A\subset S$ measurable.
\end{itemize}
We can see $\Pi$ as an element of the measurable space $(\Omega_S, \mathcal F_S)$ where $\Omega_S$ is the set of countable subsets of $S$ and $\mathcal F_S$ is the smallest $\sigma$-algebra for which the map $\Pi \mapsto \#\{\Pi \cap\nobreak B\}$ is measurable for all $B \in \mathcal S$. If~$\mu$ has no atoms, meaning no singleton with positive probability, and is $\sigma$-finite, meaning that $S$ is a countable union of measurable sets with finite measure, then a Poisson process with mean measure $\mu$ always exists (see \cite{Kingman93}, the existence theorem in section 2.5 for a proof). If~$\mu$ is absolutely continuous with respect to the Lebesgue measure, then the function \hbox{$\lambda \colon S \to \RR_{\ge 0}$} such that $\mu(A) = \riemint{A}{}{\lambda(x)}{x}$ for all measurable subsets $A \subseteq S$ is called the \textit{intensity} of the Poisson process. Here are a few important propositions about Poisson processes that we will use in the paper and proved in \cite{Kingman93}:

\begin{prop}[Mapping theorem, {\cite[\S2.3]{Kingman93}}]
\label{prel:map}
If $\Pi$ is a Poisson process with mean measure $\mu$ on $S$, and $f:S \to T$ is a measurable function such that $\mu^*(B) \coloneqq \mu(f^{-1}(B))$ has no atoms, then $f(\Pi)$ is a Poisson process on $T$ with measure $\mu^*$.
\end{prop}

\begin{prop}[Coloring theorem, {\cite[\S5.1]{Kingman93}}]
\label{prel:color}
If $\Pi$ is a Poisson process with mean measure $\mu$ on $S$, and the points are randomly colored with $k$ colors such that the probability of a point receiving the color $i$ is $p_i$, and such that the color of a point is independent of different points and of the position of the point. Let $\Pi_i$ be the subset of~$\Pi$ with color $i$. Then all the $\Pi_i$ are independent Poisson processes with mean measures $\mu_i = p_i\mu$.
\end{prop}

\begin{prop}[Campbell's theorem, {\cite[\S3.2]{Kingman93}}]
\label{prel:campbell}
Let $\Pi$ be any Poisson process on $S$ with mean measure $\mu$. Let $f\colon S \to \RR$ be a measurable function. Then
\[
\Sigma = \sum_{X \in \Pi} f(X)
\]
is absolutely convergent almost surely if, and only if,
\[
\riemint{S}{}{\min\{f, 1\}}{\mu} < \infty.
\]
If this condition holds, then
\[
\EE[e^{s\Sigma}] = \exp\left(\riemint{S}{}{(e^{sf} - 1)}{\mu}\right)
\]
for any complex $s$ for which the integral converges. Moreover,
\[
\EE[\Sigma] = \riemint{S}{}{f}{\mu}
\]
if the integral converges. In~the case where it converges, we~also have
\[
\vard[\Sigma] = \riemint{S}{}{f^2}{\mu}.
\]
\end{prop}

Many probabilities or expectations that involve Poisson processes in this paper can be reformulated as
\[
\EE \sum_{X \in \Pi} f(\Pi\setminus \{X\}, X).
\]
In these cases, there is a generalization of the formula for $\EE[\Sigma]$ in Campbell's theorem, called the Mecke equation, allowing us to compute these objects:

\begin{prop}[Mecke equation, {\cite[Th.\,4.5]{LastPenrose18}}]
\label{prel:mecke}
Let $\Pi$ be a Poisson process on $S$ with a $\sigma$-finite mean measure $\mu$, and let $f: \Omega_S \times S^k \to [0, \infty)$ be measurable. Then we have
\begin{multline*}
\EE \sum_{\substack{X_1, \ldots, X_k \\ \text{all distinct}}} f(\Pi\setminus \{X_1, \ldots, X_k\}; X_1, \ldots, X_k) \\[-10pt]
= \int_S \cdots \int_S \EE[f(\Pi; x_1, \ldots, x_k)] \, \mathrm d\mu(x_1)\, \cdots\, \mathrm d\mu(x_k).
\end{multline*}
\end{prop}

\begin{lem}[Arratia \cite{Arratia02}]
\label{lem:probaj}
For $x \ge 2$ and $1 \le j \le x$, we~have
\[
\PP\bg[J_x^*=j\bg] = \frac{1}{j\log x}\bigl(1+\bigoo{\sfrac{1}{\log x}} \bigr).
\]
\end{lem}

\begin{proof}
Let $t>0$ and let $q$ be a prime power, and consider the random variable $N_q(t)$ that counts the number of points $(T^*_i, Q^*_i)$ with $T^*_i > t$ and $Q^*_i = q$. Note that $N_q(t)$ has Poisson distribution with parameter $\sfrac{\Lambda(q)}{q^{1+t}\log q}$. Moreover, let $I^*_t \coloneqq \prod q^{N_q(t)}$ with the product being over all prime powers $q$. On one hand, when~$t$ is fixed, the expectation $\EE[(I^*_t)^{-s}]$ can be expanded as a Dirichlet series whose $i$\textsuperscript{th} coefficient equals $\PP[I^*_t = i]$. This Dirichlet series converges absolutely for $\Re(s)\ge0$. On the other hand, we~have
\[
\EE[(I^*_t)^{-s}] = \prod_q \EE\bg[q^{-sN_q(t)}\bg] = \prod_q \exp\!\bgg(\tfrac{\Lambda(q)}{q^{1+t}\log q} (q^{-s}-1)\bgg) = \frac{\zeta(1+t+s)}{\zeta(1+t)}.
\]
It follows that $\PP[I^*_t = i] = \sfrac{i^{-1-t}}{\zeta(1+t)}$ for all $i \ge 1$. All points of $\RS^*$ have distinct $T^*$-coordinates with probability one. In~this situation, we~have that $J_x^* = j$ if, and only if, there exists exactly one point $(T^*, Q^*) \in \RS^*$ such that $I^*_{T^*} = j$ and $Q^* > \sfrac{x}{j}$. Thus, we~have
\[
\one _{J_x^* = j} = \sum_{(T^*, Q^*) \in \RS^*} \one _{I^*_{T^*} = j \text{ and } Q^{*}> x/j}
\]
almost surely. Taking expectations on both sides and using the Mecke equation (Proposition \ref{prel:mecke}), we~obtain the distribution of $J_x^*$:
\begin{equation}
\label{eq:probaJ}
\PP\bg[J_x^*=j\bg]
= \riemint{0}{\infty}{\bigg(\sum_{q > x/j} \frac{\Lambda(q)}{q^{1+t}}\bigg)\cdot \frac{j^{-1-t}}{\zeta(1+t)}}{t}.
\end{equation}
We want to estimate this integral. Let $S(u) \coloneqq \sum_{q \le u} \sfrac{\Lambda(q)}{q}$. We~have $S(u) = \log u + O(1)$ for $u \ge 1$ by Mertens's estimate \cite[Th.\,3.4(a)]{Koukoulo19}. We~use partial summation to~get
\[
\sum_{q > x/j} \frac{\Lambda(q)}{q^{1+t}} = \riemint{x/j}{\infty}{u^{-t}}{S(u)} = \frac{(x/j)^{-t}}{t}\,(1+O(t)).
\]
for all $t > 0$. By putting this estimate in \eqref{eq:probaJ}, we~have
\[
\PP[J_x^* = j] = \frac{1}{j} \riemint{0}{\infty}{\frac{x^{-t}(1+O(t))}{t\cdot \zeta(1+t)}}{t}.
\]
Since $\zeta(1+t)\ge1$ for all $t>0$, the portion of the integral over $t\ge1$ is $\ll 1/(x\log x)$. On the other hand, if $t\in(0,1]$, we~have $1/\zeta(1+t)=t+O(t^2)$. We~conclude that
\[
\PP[J_x^* = j] = \frac{1}{j} \int_0^1 x^{-t}(1+O(t)) \,\mathrm dt + O(\sfrac{1}{jx\log x}) = \frac{1}{j\log x}\bg(1+O(\sfrac{1}{\log x})\bg).
\]
This concludes the proof.
\end{proof}

\begin{thebibliography}{99}
\bibitem{Arratia98} R. Arratia - “On the central role of scale invariant Poisson processes on (0,∞)”, in Microsurveys in discrete probability (Princeton, NJ, 1997), DIMACS Ser. Discrete Math. Theoret. Comput. Sci., vol. 41, American Mathematical Society, Providence, RI, 1998, p. 21-41
\bibitem{Arratia02} R. Arratia - “On the amount of dependence in the prime factorization of a uniform random integer”, in Contemporary combinatorics, Bolyai Soc. Math. Stud., vol. 10, János Bolyai Math. Soc., Budapest, 2002, p. 29-91
\bibitem{ABT06} R. Arratia, A. D. Barbour \& S. Tavaré - “A tale of three couplings: Poisson-Dirichlet and GEM approximations for random permutations”, Combin. Probab. Comput. 15 (2006) no. 1-2, p. 31-62
\bibitem{BareikisManstavicius07} G. Bareikis \& E. Manstavičius - “On the DDT theorem”, Acta Arith. 126 (2007) no. 2, p. 155-168
\bibitem{Bill72} P. Billingsley - “On the distribution of large prime divisors”, Period. Math. Hungar. 2 (1972), p. 283-289
\bibitem{delaBretecheTenenbaum16} R. de la Bretèche \& G. Tenenbaum - “Sur les processus arithmétiques liés aux diviseurs”, Adv. in Appl. Probab. 48 (2016) no. A, p. 63-76
\bibitem{DDT79} J.-M. Deshouillers, F. Dress \& G. Tenenbaum - “Lois de répartition des diviseurs. I”, Acta Arith. 34 (1979) no. 4, p. 273-285 (loose errata)
\bibitem{DonnellyTavare87} P. Donnelly \& S. Tavaré - “The population genealogy of the infinitely-many neutral alleles model”, J. Math. Biol. 25 (1987) no. 4, p. 381-391
\bibitem{Feng10} S. Feng - The Poisson-Dirichlet distribution and related topics. Models and asymptotic behaviors, Probability and its Appl., Springer, Heidelberg, 2010
\bibitem{JKB} N. L. Johnson, S. Kotz \& N. Balakrishnan - Discrete multivariate distributions, Wiley Series in Probab. and Stat.: Applied Probab. and Stat., John Wiley \& Sons, Inc., New York, 1997
\bibitem{Kingman93} J. F. C. Kingman - Poisson processes, Oxford Studies in Probab., vol. 3, The Clarendon Press, Oxford University Press, New York, 1993
\bibitem{KotzBal00} S. Kotz, N. Balakrishnan \& N. L. Johnson - Continuous multivariate distributions. Vol. 1. Models and applications, Wiley Series in Probab. and Stat.: Applied Probab. and Stat., Wiley-Interscience, New York, 2000
\bibitem{Koukoulo19} D. Koukoulopoulos - The distribution of prime numbers, Graduate Studies in Math., vol. 203, American Mathematical Society, Providence, RI, 2019
\bibitem{LastPenrose18} G. Last \& M. Penrose - Lectures on the Poisson process, Institute of Math. Stat. Textbooks, vol. 7, Cambridge University Press, Cambridge, 2018
\bibitem{Ken22} S.-K. Leung - “Dirichlet law for factorisation of integers, polynomials and permutations”, Math. Proc. Cambridge Philos. Soc. 175 (2023) no. 3, p. 649-676
\bibitem{KenWeb24} S.-K. Leung - “Dirichlet law with arbitrary parameters for factorizations of integers” (2025), Retrieved on 24 August 2025 from https://sites.google.com/view/sunkaileung/miscellaneous
\bibitem{MV07} H. L. Montgomery \& R. C. Vaughan - Multiplicative number theory. I. Classical theory, Cambridge Studies in Advanced Math., vol. 97, Cambridge University Press, Cambridge, 2007
\bibitem{NyandwiSmati13} S. Nyandwi \& A. Smati - “Distribution laws of pairs of divisors”, Integers 13 (2013), article ID A13, 13 pages
\bibitem{Tenenbaum00} G. Tenenbaum - “A rate estimate in Billingsley’s theorem for the size distribution of large prime factors”, Q. J. Math. 51 (2000) no. 3, p. 385-403
\end{thebibliography}
\end{document}
